% ============================================================================
% Fichier  : partie2/P02_C09_decomposer_espace_audit.tex
% Titre    : Décomposer l'Espace d'Audit : Gabarit Canonique et Fiche d'Assertion
% Auteur   : MINKA MI NGUIDJOÏ Thierry Emmanuel
% Version  : 1.0 -- novembre 2026
% ============================================================================
% Base : ACSI C07 (gabarit à huit niveaux d'A0 ; fiche d'assertion ; D.4.2 ;
% atomicité ; domaine D 5 × 5 ; lecture D.1 4/0/0, D.2 1/1/2, D.3 4/0/0,
% D.4 2/0/2, D.5 2/1/1).
% Sources : A0 (hal-05701007), thématique T2 et gabarit canonique ; Miller
% (1956) pour la borne cognitive ; Simon (1962) pour la décomposabilité des
% systèmes complexes ; Minto (2009) pour le principe « mutuellement
% exclusif, collectivement exhaustif ». Décision du 06/11/2026 : aucun
% renvoi à un article non publié (A4) ; les conditions sont fondées sur ces
% travaux publiés.
% Classeur O4 : dix onglets (NOTICE, MISSION, PARAMÈTRES, DOMAINES, GABARIT
% A4, RÉPONDANTS, MESURES, CYCLES, CONTRADICTOIRE, SYNTHÈSE).
% Corrigés : corriges/Corr_C09.tex.
% ============================================================================

\chapter[Décomposer l'Espace d'Audit]{Décomposer l'Espace d'Audit : Gabarit Canonique et Fiche d'Assertion}
\label{chap:decomposition}
\chaptermark{Décomposer l'espace d'audit}

\epigraph{%
    « Ce qu'on ne peut pas tenir en tête,\\
    on ne peut pas l'auditer~: on le découpe. »%
}{Principe de la méthode}

\niveauChapitre{Intermédiaire}{%
    Maîtrise des Chapitres~\ref{chap:controle-interne} (contrôles)
    et~\ref{chap:planifier} (plan d'audit)%
}

\begin{boxobjectifs}
\begin{itemize}
    \item Justifier la décomposition de l'espace d'audit par les limites de
          la mémoire de travail et la structure des systèmes complexes.
    \item Décomposer un domaine selon le gabarit canonique à huit niveaux~:
          domaine, sous-domaine, processus, activité, contrôle, test, preuve,
          trace.
    \item Rédiger la fiche d'assertion d'un contrôle à plusieurs critères, et
          en déduire la cote du contrôle.
    \item Appliquer la règle d'atomicité aux objectifs, aux critères et aux
          questions.
    \item Vérifier qu'une décomposition est exhaustive, non redondante et
          respecte la borne cognitive, à la main et par un programme.
    \item Construire la décomposition dans le classeur O4.
    \item Reconnaître les domaines où le gabarit doit être adapté.
\end{itemize}
\end{boxobjectifs}

\bigskip

% ============================================================================
\section*{Question directrice : pourquoi la grille de 45 questions a-t-elle échoué~?}
% ============================================================================

Le prologue a donné deux raisons à l'échec de la première mission de
LiZbiZ~: une grille binaire qui interdisait de dire « je ne sais pas », et
une grille reprise d'une autre entreprise, sans critères propres. Il y en a
une troisième. Les neuf questions du domaine~D mélangeaient des objets de
niveaux différents~: une question portait sur la politique d'accès, une
autre sur un journal technique, une troisième combinait deux conditions.
Personne ne pouvait dire si le domaine avait été couvert en entier.

La seconde mission a commencé par \emph{décomposer} le domaine~D en
$5 \times 5$ éléments~\cite{minka2026auditologie}. Ce chapitre enseigne
cette opération.

\medskip
\textbf{Plan du chapitre.}
\begin{enumerate}
    \item Pourquoi décomposer (section~\ref{sec:c09-pourquoi}).
    \item Le gabarit canonique (section~\ref{sec:c09-gabarit}).
    \item La fiche d'assertion (section~\ref{sec:c09-fiche}).
    \item Les règles de qualité (section~\ref{sec:c09-regles}).
    \item La décomposition dans le classeur O4
          (section~\ref{sec:c09-o4}).
    \item Le domaine D de LiZbiZ (section~\ref{sec:c09-lizbiz}).
    \item Les limites du gabarit (section~\ref{sec:c09-limites}).
\end{enumerate}

\separateur

% ============================================================================
\section{Pourquoi Décomposer}
\label{sec:c09-pourquoi}
% ============================================================================

\subsection{La borne de la mémoire de travail}

Un auditeur ne peut pas examiner « le contrôle d'accès » d'un bloc. Les
travaux de George Miller ont établi que la mémoire de travail ne retient
qu'un petit nombre d'éléments à la fois, de l'ordre de sept, plus ou moins
deux~\cite{miller1956}. L'auditologie en fait l'un de ses axiomes
(axiome~2, \emph{cognition limitée}) et en déduit qu'un instrument qui
couvre un espace d'audit non trivial doit être organisé en une
décomposition hiérarchique, dont chaque niveau respecte cette
borne~\cite{minka2026auditologie}.

\subsection{Les systèmes complexes se décomposent}

Herbert Simon a montré que les systèmes complexes, naturels ou artificiels,
sont le plus souvent hiérarchiques et « presque décomposables »~: les
interactions à l'intérieur d'un sous-système sont plus fortes que les
interactions entre sous-systèmes~\cite{simon1962complexity}. Un système
d'information en est un exemple~: la gestion des comptes interagit
fortement avec l'authentification, faiblement avec la gestion des
sauvegardes. C'est ce qui permet d'auditer chaque sous-domaine
séparément, puis d'en assembler les conclusions.

\begin{boxdef}
\termecle{Espace d'audit}\index{espace d'audit}~: ensemble des objets sur
lesquels porte une mission (processus, contrôles, données, traces), tel que
le délimite la lettre de mission.

\medskip
\termecle{Décomposition}\index{décomposition}~: découpage hiérarchique de
l'espace d'audit en niveaux emboîtés, chaque élément d'un niveau se divisant
en éléments du niveau inférieur.
\end{boxdef}

\subsection{Ce que la décomposition apporte}

\begin{enumerate}
    \item \textbf{La couverture}~: on sait ce qui a été audité, et ce qui ne
          l'a pas été.
    \item \textbf{La traçabilité}~: chaque conclusion remonte jusqu'à une
          trace dans le système, ce qui la rend opposable (axiome~6).
    \item \textbf{La localisation de l'incertitude}~: les réponses \NEK{} se
          situent à un endroit précis de l'arbre, et non dans « le domaine »
          (section~\ref{sec:c09-lizbiz}).
    \item \textbf{La répartition du travail}~: chaque sous-domaine peut être
          confié à un auditeur.
\end{enumerate}

% ============================================================================
\section{Le Gabarit Canonique}
\label{sec:c09-gabarit}
% ============================================================================

\subsection{Huit niveaux}

Le programme de l'auditologie propose un \termecle{gabarit
canonique}\index{gabarit canonique} de décomposition à huit
niveaux~\cite{minka2026auditologie}.

\noindent\begin{minipage}{\linewidth}
\begin{center}
\begin{tikzpicture}[
    niv/.style = {
        draw=CouleurBordDef, fill=CouleurDef, rounded corners=3pt,
        minimum width=3.1cm, minimum height=0.62cm, align=center,
        font=\footnotesize\bfseries\color{CouleurPrimaire}
    },
    note/.style = {font=\scriptsize\itshape\color{CouleurSecondaire}, anchor=west, text width=6.4cm},
    fleche/.style = {-{Stealth[length=4pt]}, thick, color=CouleurSecondaire}
]
\foreach \t/\q [count=\i] in {
    1. Domaine/Que soumet-on à l'audit~?,
    2. Sous-domaine/Quels grands ensembles le composent~?,
    3. Processus/Quelles suites d'activités~?,
    4. Activité/Quelles actions élémentaires~?,
    5. Contrôle/Quel dispositif maîtrise l'activité~?,
    6. Test/Comment l'auditeur le vérifie-t-il~?,
    7. Preuve/Quel élément le test produit-il~?,
    8. Trace/Quel enregistrement du système l'atteste~?} {
    \node[niv] (n\i) at (0, -0.85*\i) {\t};
    \node[note] at (1.8, -0.85*\i) {\q};
}
\foreach \i [evaluate=\i as \j using int(\i+1)] in {1,...,7}
    \draw[fleche] (n\i) -- (n\j);
\draw[decorate, decoration={brace, amplitude=5pt, mirror}, thick, color=CouleurAccent]
    (-1.75,-0.55) -- (-1.75,-3.65) node[midway, left=6pt, font=\scriptsize, align=right,
    color=CouleurAccent] {ce qui\\ est audité};
\draw[decorate, decoration={brace, amplitude=5pt, mirror}, thick, color=CouleurAccent]
    (-1.75,-3.95) -- (-1.75,-7.05) node[midway, left=6pt, font=\scriptsize, align=right,
    color=CouleurAccent] {comment on\\ le vérifie};
\end{tikzpicture}
\end{center}
\vspace{-0.4em}
\begin{center}{\small\itshape Figure~9.1. Le gabarit canonique de
décomposition (d'après~\cite{minka2026auditologie})}\end{center}
\end{minipage}

\medskip

Les quatre premiers niveaux décrivent \emph{ce qui est audité}~: ils
découpent le système. Les quatre derniers décrivent \emph{comment on le
vérifie}~: du dispositif de maîtrise jusqu'à l'enregistrement qui en prouve
le fonctionnement. Le contrôle fait la charnière~: c'est l'objet sur lequel
l'auditeur se prononce.

\subsection{Les huit niveaux en détail}

\begin{tableaucours}{L{2.3cm}L{5.0cm}L{4.9cm}}{%
    \tableauHeader{Niveau} &
    \tableauHeader{Définition} &
    \tableauHeader{Exemple (domaine D de LiZbiZ)}
}
Domaine & Champ d'audit défini par la lettre de mission & Contrôle d'accès
à l'application intégrée \\
Sous-domaine & Grand ensemble cohérent du domaine & Cycle de vie des
comptes \\
Processus & Suite d'activités qui produit un résultat & Départ d'un
agent \\
Activité & Action élémentaire du processus & Retrait des accès \\
Contrôle & Dispositif qui maîtrise l'activité & Désactivation du compte
sous 48~heures \\
Test & Procédure par laquelle l'auditeur vérifie le contrôle & Rapprocher
les départs et les comptes actifs \\
Preuve & Élément produit et conservé par le test & Extraction des comptes
actifs, avec son empreinte \\
Trace & Enregistrement du système qui atteste le fonctionnement &
Entrée du journal de désactivation \\
\end{tableaucours}
\vspace{-0.8em}
\begin{center}{\small\itshape Tableau~9.1. Les niveaux du gabarit
canonique}\end{center}

\begin{boxalert}
\textbf{Violation typique.} Placer une politique (« une politique d'accès
approuvée existe ») au niveau du test, et un journal technique au niveau du
sous-domaine, dans la même liste.

\smallskip
\textbf{Règle~:} chaque élément se place à un seul niveau, celui qui
correspond à sa nature. Un mélange de niveaux rend impossible le contrôle de
la couverture.
\end{boxalert}

\subsection{Preuve et trace~: deux choses différentes}

La \emph{trace} existe dans le système, que l'auditeur passe ou non~: une
entrée de journal, un champ de date de désactivation. La \emph{preuve}
est ce que l'auditeur en tire et conserve~: l'extraction horodatée, son
empreinte, la requête qui l'a produite. Une trace qui n'est pas extraite
n'est pas une preuve~; une preuve qui ne renvoie à aucune trace (une
capture d'écran sans contexte) est fragile. Le Chapitre~\ref{chap:preuve}
en tirera la hiérarchie de fiabilité.

% ============================================================================
\section{La Fiche d'Assertion}
\label{sec:c09-fiche}
% ============================================================================

\subsection{Cinq attributs par contrôle}

Le gabarit dit \emph{où} se trouve un contrôle. La \termecle{fiche
d'assertion}\index{fiche d'assertion} dit \emph{ce que l'auditeur affirme}
à son sujet, et ce qui arrive si l'affirmation est fausse.

\begin{tableaucours}{L{2.4cm}L{5.0cm}L{4.8cm}}{%
    \tableauHeader{Attribut} &
    \tableauHeader{Rôle} &
    \tableauHeader{Forme}
}
Objectif & Ce que le contrôle doit garantir & À l'infinitif~: « S'assurer
que\ldots » \\
Critère & La condition vérifiable qui définit la conformité & Au présent~:
une condition unique, observable \\
Question & La question posée pour coter le contrôle & Interrogative~;
réponse \OUI{}, \NON{} ou \NEK{} \\
Risque & L'état de non-conformité si la réponse est \NON{} & L'antithèse
du critère \\
Conséquence & L'effet sur l'organisation si le risque se réalise &
Financier, opérationnel, réglementaire \\
\end{tableaucours}
\vspace{-0.8em}
\begin{center}{\small\itshape Tableau~9.2. Les attributs de la fiche
d'assertion}\end{center}

La fiche est complétée par les niveaux inférieurs du gabarit~: le test,
la preuve attendue et la trace. Le prologue en a donné un premier exemple
(contrôle D.2.1)~; en voici un second.

\begin{boxobjet}{D.4.2. Revue des comptes à privilèges}
\textbf{Position~:} domaine D $\rightarrow$ sous-domaine D.4 (comptes à
privilèges) $\rightarrow$ processus «~gestion des comptes d'administration~»
$\rightarrow$ activité «~revue des comptes~» $\rightarrow$ contrôle D.4.2.

\textbf{Objectif~:} s'assurer que chaque compte d'administration de
l'application correspond à un besoin en cours.

\textbf{Critère~1~:} la liste des comptes d'administration est revue chaque
trimestre par le responsable de la sécurité.

\textbf{Question~1~:} la revue des comptes d'administration a-t-elle été
réalisée pour chaque trimestre de la période~?

\textbf{Critère~2~:} le compte rendu de chaque revue est signé par le
responsable de la sécurité.

\textbf{Question~2~:} le compte rendu de chaque revue trimestrielle est-il
signé par le responsable de la sécurité~?

\textbf{Test~:} demander les comptes rendus des deux trimestres de la
période~; comparer leur liste à l'extraction des comptes d'administration.

\textbf{Preuve et trace~:} comptes rendus signés~; extraction des comptes
d'administration, avec son empreinte~; dates de dernière connexion de ces
comptes dans le journal.

\textbf{Risque~:} un compte d'administration inutile reste actif et sert à
modifier des données ou des droits sans contrôle.

\textbf{Conséquence~:} fraude ou altération non détectée, perte de fiabilité
des états de stock, manquement à l'obligation de sécuriser le système
(loi n°~2010/012, art.~26).
\end{boxobjet}

\subsection{De la fiche à la cotation}

La question de la fiche est celle que l'auditeur cote \OUI{}, \NON{} ou
\NEK{} (Chapitre~\ref{chap:preuve}). Le risque est \emph{activé} lorsque la
réponse est \NON{}~; il reste \emph{indéterminé} lorsque la réponse est
\NEK{}, et c'est précisément ce que mesure le taux d'incertitude~$U$.

Un contrôle peut porter plusieurs critères, comme D.4.2~: une revue
réalisée mais non signée ne satisfait pas le contrôle, et l'auditeur doit
pouvoir dire lequel des deux fait défaut. Chaque critère reçoit sa propre
question et sa propre cote. La cote du contrôle s'en déduit.

\begin{boxinfo}
\textbf{Règle de composition des cotes.} Un seul critère \NON{} suffit pour
que le contrôle soit \NON{}~; à défaut, un seul critère \NEK{} donne
\NEK{}~; le contrôle n'est coté \OUI{} que si tous ses critères le sont.
En notation~: si $c_1, \dots, c_k$ sont les cotes des critères,
\[
\text{cote} =
\begin{cases}
\NON{} & \text{si } \exists\, i,\ c_i = \NON{},\\
\NEK{} & \text{sinon, si } \exists\, i,\ c_i = \NEK{},\\
\OUI{} & \text{sinon.}
\end{cases}
\]
\end{boxinfo}

\begin{tableaucours}{L{3.2cm}L{3.2cm}cL{4.2cm}}{%
    \tableauHeader{Critère 1 (revue faite)} &
    \tableauHeader{Critère 2 (signée)} &
    \tableauHeader{Cote D.4.2} &
    \tableauHeader{Lecture}
}
\OUI{} & \OUI{} & \OUI{} & Contrôle conforme \\
\OUI{} & \NON{} & \NON{} & Revue faite, non signée \\
\NEK{} & \OUI{} & \NEK{} & Signée, mais on ignore si elle couvre la
période \\
\NON{} & \NEK{} & \NON{} & Un défaut établi suffit \\
\end{tableaucours}
\vspace{-0.8em}
\begin{center}{\small\itshape Tableau~9.3. Composer les cotes de deux
critères}\end{center}

% ============================================================================
\section{Les Règles de Qualité}
\label{sec:c09-regles}
% ============================================================================

\subsection{L'atomicité}

\begin{boxinfo}
\textbf{Règle d'atomicité.} L'objectif, le critère et la question ne
contiennent aucun connecteur logique «~et~» ni «~ou~». Toute assertion qui
en contient est scindée en autant d'assertions qu'il y a de conditions
indépendantes.
\end{boxinfo}

L'atomicité garantit qu'une réponse \NON{} désigne un seul défaut. Une
question double, «~les comptes des agents partis sont-ils désactivés
\emph{et} leurs données archivées~?~», n'admet pas de réponse honnête
lorsque les comptes sont désactivés mais les données non archivées~: le
répondant choisit au hasard, ou répond \NEK{} à tort.

\begin{boxalert}
\textbf{Violation typique.} « Les mots de passe sont-ils robustes et
changés régulièrement~? »

\smallskip
\textbf{Règle~:} scinder~: « La politique impose-t-elle au moins douze
caractères~? » et « Les mots de passe des comptes d'administration ont-ils
été changés au cours des 90 derniers jours~? ». Chaque question a son
critère, sa preuve et sa trace.
\end{boxalert}

\subsection{Exhaustivité, non-redondance, borne cognitive}

Une décomposition est de qualité lorsqu'elle respecte trois conditions. Les
deux premières sont connues des consultants sous le nom de principe
« mutuellement exclusif, collectivement exhaustif »~\cite{minto2009pyramid}~;
la troisième vient de la borne de Miller.

\begin{tableaucours}{L{3.2cm}L{5.0cm}L{4.0cm}}{%
    \tableauHeader{Condition} &
    \tableauHeader{Sens} &
    \tableauHeader{Défaut si elle manque}
}
Exhaustivité & Les éléments d'un niveau couvrent entièrement leur parent &
Une partie du domaine n'est pas auditée \\
Non-redondance & Deux éléments d'un même niveau ne se recouvrent pas~;
chaque élément dérive d'un seul parent & Double comptage, conclusions
contradictoires \\
Borne cognitive & Aucun élément n'a plus de sept enfants environ & Grille
illisible, éléments oubliés \\
\end{tableaucours}
\vspace{-0.8em}
\begin{center}{\small\itshape Tableau~9.4. Les trois conditions d'une bonne
décomposition}\end{center}

\subsection{Vérifier par programme}

Deux des trois conditions se vérifient automatiquement sur la liste des
codes du classeur. L'exhaustivité, elle, demande un jugement~: on compare
les sous-domaines au périmètre de la lettre de mission et aux référentiels.

\begin{boxoutils}{Contrôler un arbre de décomposition}
\begin{lstlisting}[language=Python]
from collections import Counter
codes = ["D", "D.1", "D.2", "D.3", "D.4", "D.5",
         "D.2.1", "D.2.2", "D.2.3", "D.2.4", "D.2.5"]   # extrait
parent = lambda c: c.rsplit(".", 1)[0] if "." in c else None
# 1. chaque code a un parent présent (pas d'élément orphelin)
orphelins = [c for c in codes if parent(c) and parent(c) not in codes]
# 2. aucun code en double (non-redondance syntaxique)
doublons = [c for c, n in Counter(codes).items() if n > 1]
# 3. borne cognitive : au plus 7 enfants par élément
enfants = Counter(parent(c) for c in codes if parent(c))
trop = {p: n for p, n in enfants.items() if n > 7}
print(orphelins, doublons, trop)        # attendu : [] [] {}
\end{lstlisting}
\end{boxoutils}

\begin{boxprocedure}
\textbf{La décomposition revue par le chef de mission.} Pour que la
décomposition puisse être revue, le classeur indique, pour chaque élément~:
\begin{itemize}
    \item son code (par exemple D.4.2), son niveau et son parent~;
    \item pour chaque contrôle, sa fiche d'assertion complète, critère par
          critère~;
    \item pour chaque test, la preuve attendue et la référence du papier de
          travail où elle sera classée~;
    \item le référentiel d'où provient chaque critère (procédure interne,
          COBIT, ISO/IEC~27002, loi)~;
    \item la date de la revue de la décomposition et le nom du réviseur.
\end{itemize}
\end{boxprocedure}

\subsection{Relier la décomposition aux référentiels}

Les référentiels de contrôle fournissent des critères~: la
mesure~5.18 de l'ISO/IEC~27002:2022 sur les droits d'accès, les pratiques de
l'objectif DSS05 de COBIT, les principes du COSO
(Chapitre~\ref{chap:referentiels}). Ils ne fournissent pas la
décomposition de \emph{votre} mission~: c'est l'auditeur qui la construit, à
partir de la lettre de mission et de la prise de connaissance, puis va
chercher dans les référentiels les critères de chaque contrôle.

% ============================================================================
\section{La Décomposition dans le Classeur O4}
\label{sec:c09-o4}
% ============================================================================

Le classeur d'aide à l'audit O4 porte toute la mission. Il compte dix
onglets~; trois d'entre eux servent à la décomposition.

\begin{tableaucours}{L{3.2cm}L{9.0cm}}{%
    \tableauHeader{Onglet} &
    \tableauHeader{Contenu}
}
NOTICE & Mode d'emploi, conventions, règle de cotation \\
MISSION & Référence de la mission, objectifs, période, équipe \\
PARAMÈTRES & Seuils de lecture ($\mu$, $U$), listes de valeurs \\
DOMAINES & Liste des domaines et sous-domaines, avec leur responsable \\
GABARIT A4 & Une ligne par contrôle, sur les huit niveaux~; la fiche
d'assertion est repliée sous la ligne, un critère par sous-ligne, avec sa
source, sa question et sa cote \\
RÉPONDANTS & Personnes interrogées, fonction, date \\
MESURES & Calcul automatique de $\mu$, $U$, $H$ par sous-domaine et par
domaine \\
CYCLES & Comparaison avec la mission précédente (Chapitre~\ref{chap:suivi}) \\
CONTRADICTOIRE & Position de l'audité et position finale, par constat \\
SYNTHÈSE & Lecture par domaine, pour le rapport \\
\end{tableaucours}
\vspace{-0.8em}
\begin{center}{\small\itshape Tableau~9.5. Les dix onglets du classeur
O4}\end{center}

\begin{boxprocedure}
\textbf{Construire un domaine dans O4, pas à pas}
\begin{enumerate}
    \item Dans DOMAINES, créer le domaine (D) et ses sous-domaines (D.1 à
          D.5), avec un auditeur responsable de chacun.
    \item Dans GABARIT A4, créer une ligne par contrôle, en renseignant les
          huit niveaux (le code du contrôle reprend celui du processus~:
          D.2.1 est le premier contrôle de D.2).
    \item Sous chaque contrôle, déplier la fiche d'assertion~: une
          sous-ligne par critère, avec la source et la question.
    \item Laisser vide la colonne des cotes~: elle se remplit en phase~2.
    \item Faire revoir l'onglet par le chef de mission, et inscrire la date
          de revue dans MISSION.
\end{enumerate}
\end{boxprocedure}

% ============================================================================
\section{Application : le Domaine D de LiZbiZ}
\label{sec:c09-lizbiz}
% ============================================================================

\subsection{Cinq sous-domaines}

\begin{boxscenario}{Opération LIZBIZ. Décomposer le domaine D}
AUDIX CONSEIL décompose le domaine~D (contrôle d'accès à l'application
intégrée) en cinq sous-domaines, eux-mêmes découpés en cinq éléments au
plus, soit $5 \times 5$ éléments.

\medskip
\begin{tableaucours}{cL{3.6cm}L{6.8cm}}{%
    \tableauHeader{Code} &
    \tableauHeader{Sous-domaine} &
    \tableauHeader{Processus}
}
D.1 & Politique et organisation & Politique d'accès~; rôles~; profils
types~; validation des demandes~; communication \\
D.2 & Cycle de vie des comptes & Arrivée~; changement de fonction~; départ~;
comptes génériques~; revue périodique \\
D.3 & Authentification & Mots de passe~; double authentification~;
verrouillage~; sessions~; comptes par défaut \\
D.4 & Comptes à privilèges & Attribution~; revue~; usage tracé~; comptes de
service~; secours \\
D.5 & Journalisation et revue & Journaux d'accès~; protection des
journaux~; revue des anomalies~; conservation~; alertes \\
\end{tableaucours}

\medskip
\textbf{Vérification.} Chaque sous-domaine compte cinq processus, sous la
borne de sept. Les cinq sous-domaines couvrent l'ensemble du contrôle
d'accès décrit dans la lettre de mission. Le compte générique de
l'agence de Garoua relève de D.2 et non de D.4~: il n'a pas de privilèges
d'administration.
\end{boxscenario}

\subsection{Localiser l'incertitude}

\begin{boxlecture}{Domaine D décomposé}
Le prologue a lu le domaine~D d'un bloc~: 20 réponses, $\mu = 0{,}65$,
$U = 0{,}25$, $H \approx 1{,}24$~bit. La décomposition permet de répartir ces
réponses~: une question par sous-domaine, posée aux quatre mêmes
interlocuteurs.

\medskip
\begin{tableaucours}{cL{3.6cm}cccc}{%
    \tableauHeader{Code} &
    \tableauHeader{Sous-domaine} &
    \tableauHeader{\OUI{}/\NON{}/\NEK{}} &
    \tableauHeader{$\mu$} & \tableauHeader{$U$} & \tableauHeader{$H$}
}
D.1 & Politique et organisation & 4 / 0 / 0 & 1,00 & 0,00 & 0,00 \\
D.2 & Cycle de vie des comptes & 1 / 1 / 2 & 0,25 & 0,50 & 1,50 \\
D.3 & Authentification & 4 / 0 / 0 & 1,00 & 0,00 & 0,00 \\
D.4 & Comptes à privilèges & 2 / 0 / 2 & 0,50 & 0,50 & 1,00 \\
D.5 & Journalisation et revue & 2 / 1 / 1 & 0,50 & 0,25 & 1,50 \\
\end{tableaucours}

\medskip
\textbf{Constatation.} Toute l'incertitude du domaine se concentre dans
trois sous-domaines~: D.2, D.4 et D.5. Les sous-domaines D.1 et D.3 sont
connus et conformes.

\medskip
\textbf{Analyse} (confiance moyenne). Lu d'un bloc, le domaine paraissait
uniformément opaque. Décomposé, il montre que la politique et
l'authentification sont maîtrisées, et que l'ignorance porte sur la vie des
comptes, les privilèges et les journaux~: c'est là que les tests doivent
porter. D.1 et D.3, avec une entropie nulle et $U = 0$, sont des
certitudes~; à l'inverse, une entropie nulle avec $U = 1$ serait une
ignorance unanime (Chapitre~\ref{chap:conclure}).

\medskip
\hfill\lecture{0{,}65}{0{,}25}{1{,}24}
\end{boxlecture}

\begin{boxinfo}
\textbf{Vérifier l'agrégation.} Les totaux des cinq sous-domaines
redonnent ceux du domaine~: $4+1+4+2+2 = 13$~\OUI{}, $0+1+0+0+1 = 2$~\NON{},
$0+2+0+2+1 = 5$~\NEK{}, soit 13/2/5 sur 20. L'entropie du domaine (1,24)
n'est pas la moyenne des entropies des sous-domaines (0,80)~: elle mesure
aussi la dispersion \emph{entre} sous-domaines. On recalcule toujours les
mesures d'un niveau à partir des réponses, jamais à partir des mesures du
niveau inférieur.
\end{boxinfo}

\begin{boxmetier}{Auditeur SI}
\textbf{Construire la décomposition d'un domaine}
\begin{enumerate}
    \item Partir du domaine défini par la lettre de mission, sans
          l'élargir.
    \item Découper en cinq à sept sous-domaines, à partir de la prise de
          connaissance et des référentiels.
    \item Descendre jusqu'aux contrôles, puis rédiger leurs fiches
          d'assertion, critère par critère.
    \item Pour chaque contrôle, définir le test, la preuve attendue et la
          trace à rechercher.
    \item Faire revoir la décomposition par le chef de mission avant le
          terrain, en vérifiant les trois conditions.
\end{enumerate}
\end{boxmetier}

\begin{boxmetier}{Auditeur de sécurité}
\textbf{Puiser les critères dans les référentiels techniques.} Pour D.3
(authentification), l'auditeur de sécurité tire ses critères de la
politique interne, de l'ISO/IEC~27002 (mesure 5.17, information
d'authentification) et des guides de configuration (Chapitre~\ref{chap:audit-technique}).
Il note la source de chaque critère dans O4~: sans source, un critère est
une opinion.
\end{boxmetier}

% ============================================================================
\section{Les Limites du Gabarit}
\label{sec:c09-limites}
% ============================================================================

Le gabarit n'est pas une loi universelle. Il suppose trois conditions sur le
domaine audité, dont l'échec appelle une adaptation, et non l'abandon du
gabarit.

\begin{tableaucours}{L{3.6cm}L{4.4cm}L{4.2cm}}{%
    \tableauHeader{Condition} &
    \tableauHeader{Situation où elle manque} &
    \tableauHeader{Adaptation}
}
Richesse suffisante & Un audit de conformité qui vérifie seulement
l'existence de documents & Gabarit tronqué (domaine, sous-domaine,
contrôle, test) \\
Cohérence sémantique & Un domaine nouveau, sans regroupements stables
(audit d'un système d'IA) & Décomposition provisoire, revue à chaque
mission \\
Hétérogénéité bornée & Une mission qui mêle sécurité, finance et risques
opérationnels & Plusieurs arbres, un par domaine \\
\end{tableaucours}
\vspace{-0.8em}
\begin{center}{\small\itshape Tableau~9.6. Conditions de validité et
adaptations du gabarit}\end{center}

\begin{boiteRemarque}{Et la première version de la méthode~?}
Une première version de la méthode du cours décomposait l'assertion
d'audit elle-même en huit niveaux~: domaine, point de contrôle, objectif,
critère, élément requis, question, risque, conséquence. Ces notions sont
conservées, mais à leur place~: les deux premières correspondent aux
niveaux supérieurs du gabarit, l'élément requis devient la preuve, et les
cinq autres forment la fiche d'assertion de chaque contrôle. Le gabarit
découpe le \emph{système}~; la fiche porte l'\emph{affirmation}.
\end{boiteRemarque}

\begin{boxinfo}
\textbf{Ce que dit la recherche.} Le gabarit canonique figure dans le
programme de l'auditologie (A0, thématique~T2)~\cite{minka2026auditologie}.
Les conditions de qualité de ce chapitre s'appuient sur des résultats
publiés~: la borne de Miller~\cite{miller1956}, la quasi-décomposabilité de
Simon~\cite{simon1962complexity} et le principe d'exclusivité et
d'exhaustivité~\cite{minto2009pyramid}. Leur formalisation propre à l'audit
reste un sujet de recherche ouvert.
\end{boxinfo}

\separateur

\begin{boxresume}
\textbf{Ce qu'il faut retenir de ce chapitre~:}
\begin{itemize}
    \item \textbf{On décompose parce que la mémoire de travail est
          limitée}~: environ sept éléments par niveau~; et parce que les
          systèmes complexes sont presque décomposables.
    \item \textbf{Le gabarit canonique} a huit niveaux~: domaine,
          sous-domaine, processus, activité (ce qui est audité), contrôle,
          test, preuve, trace (comment on le vérifie).
    \item \textbf{La fiche d'assertion} donne à chaque contrôle un objectif,
          un ou plusieurs critères, des questions, un risque et une
          conséquence~; la cote du contrôle se déduit de celles des
          critères (un \NON{}, sinon un \NEK{}, sinon \OUI{}).
    \item \textbf{L'atomicité} interdit «~et~» et «~ou~» dans l'objectif, le
          critère et la question.
    \item \textbf{Une bonne décomposition} est exhaustive, non redondante,
          et respecte la borne cognitive~; deux conditions se vérifient par
          programme.
    \item \textbf{Le classeur O4} porte la décomposition~; ses mesures se
          recalculent à chaque niveau à partir des réponses.
    \item \textbf{La décomposition localise l'incertitude}~: elle montre où
          porte l'ignorance, et donc où tester.
\end{itemize}
\end{boxresume}

\begin{boxinfo}
\textbf{Points de contrôle de l'élève.} Je sais~: citer les huit
niveaux~; distinguer preuve et trace~; rédiger une fiche d'assertion à deux
critères~; composer des cotes~; rendre une assertion atomique~; vérifier
une décomposition avec le programme du chapitre.
\end{boxinfo}

\bigskip

\begin{boxexo}[title={\ding{45}\quad Exercice~9.1. À quel niveau~? \niveaubase}]
Placez chaque élément dans le niveau du gabarit qui lui correspond.
\begin{enumerate}[(a)]
    \item La revue trimestrielle des droits par le responsable de la
          sécurité.
    \item Une ligne du journal indiquant la désactivation du compte
          « afouda » le 3~mai.
    \item La gestion des sauvegardes.
    \item L'extraction des comptes actifs au 30~avril, avec son empreinte.
    \item Rapprocher la liste des départs et celle des comptes actifs.
    \item Le départ d'un agent.
\end{enumerate}
\end{boxexo}

\begin{boxexo}[title={\ding{45}\quad Exercice~9.2. Rendre atomique \niveaubase}]
Réécrivez les assertions suivantes pour qu'elles respectent la règle
d'atomicité, et donnez pour chacune le critère et la preuve attendue.
\begin{enumerate}[(a)]
    \item « Les sauvegardes sont réalisées chaque nuit et testées chaque
          mois. »
    \item « Les comptes génériques sont interdits ou placés sous la
          responsabilité d'un titulaire nommé. »
    \item « S'assurer que les journaux sont conservés et protégés contre la
          modification. »
\end{enumerate}
\end{boxexo}

\begin{boxexo}[title={\ding{45}\quad Exercice~9.3. Composer des cotes \niveaubase}]
Un contrôle a trois critères. Donnez la cote du contrôle pour~: (a)
\OUI{}, \OUI{}, \NEK{}~; (b) \NEK{}, \NON{}, \OUI{}~; (c) \NEK{}, \NEK{},
\NEK{}~; (d) \OUI{}, \OUI{}, \OUI{}. Pourquoi la règle ne donne-t-elle
jamais \OUI{} lorsqu'un critère est \NEK{}~?
\end{boxexo}

\begin{boxexo}[title={\ding{45}\quad Exercice~9.4. Une fiche complète \niveaumoyen}]
Rédigez la fiche d'assertion complète du contrôle D.2.1 (désactivation des
comptes des agents partis sous 48~heures), avec deux critères, le test, la
preuve et la trace.
\end{boxexo}

\begin{boxexo}[title={\ding{45}\quad Exercice~9.5. Diagnostiquer une décomposition \niveaumoyen}]
Un auditeur propose, pour le domaine des sauvegardes~: B.1 politique~; B.2
sauvegardes quotidiennes~; B.3 journal des sauvegardes du serveur de
Douala~; B.4 restauration~; B.5 sauvegardes et restaurations testées~; B.6
« vérifier les sauvegardes ». Identifiez les défauts (niveau, redondance,
atomicité) et proposez une décomposition corrigée.
\end{boxexo}

\begin{boxexo}[title={\ding{45}\quad Exercice~9.6. Le programme de contrôle \niveaumoyen}]
Appliquez à la main le programme de la section~\ref{sec:c09-regles} à la
liste~: D, D.1, D.2, D.2.1, D.2.2, D.3.1, D.4, D.4.1, D.4.1. Qu'affiche-t-il~?
Quelle condition le programme ne peut-il pas vérifier, et pourquoi~?
\end{boxexo}

\begin{boxexo}[title={\ding{45}\quad Exercice~9.7. L'agrégation \niveaumoyen}]
Recalculez l'entropie du domaine~D à partir des totaux 13/2/5, puis la
moyenne des entropies des cinq sous-domaines. Expliquez l'écart.
\end{boxexo}

\begin{boxexo}[title={\ding{45}\quad Exercice~9.8. Construire dans O4 \niveaumoyen}]
Décrivez les lignes que vous créez dans l'onglet GABARIT A4 pour le
contrôle D.4.2, en précisant pour chacune les colonnes renseignées et
celles laissées vides avant le terrain.
\end{boxexo}

\begin{boxexo}[title={\ding{45}\quad Exercice~9.9. Décomposer le domaine des stocks \niveauexpert}]
La seconde zone d'audit de LiZbiZ est la gestion des stocks dans
l'application intégrée.
\begin{enumerate}[(a)]
    \item Proposez cinq sous-domaines et, pour l'un d'eux, ses processus et
          ses activités.
    \item Rédigez la fiche d'assertion complète de deux contrôles, dont celui
          qui porte sur la validation des ajustements de stock.
    \item Vérifiez les trois conditions de qualité de votre décomposition.
    \item Expliquez où vous placeriez le fichier tableur de suivi des stocks
          tenu à Garoua.
\end{enumerate}
\end{boxexo}

\begin{boxexo}[title={\ding{45}\quad Exercice~9.10. Adapter le gabarit \niveauexpert}]
La Chambre des comptes audite la conformité des paiements de l'ANPP au
régime financier de l'État~: pour chaque paiement, il s'agit de vérifier
l'existence de quatre pièces (engagement, liquidation, ordonnancement,
visa). Quelle adaptation du gabarit proposez-vous, et pourquoi~? Rédigez
la décomposition et la fiche d'assertion d'un contrôle.
\end{boxexo}

\bigskip

\begin{boxinfo}
\textbf{Pour aller plus loin.}
\begin{itemize}
    \item G.~Miller, \textit{The Magical Number Seven, Plus or Minus
          Two}~\cite{miller1956}~: l'article fondateur sur la mémoire de
          travail.
    \item H.~Simon, \textit{The Architecture of
          Complexity}~\cite{simon1962complexity}~: pourquoi les systèmes
          complexes se décomposent.
    \item B.~Minto, \textit{The Pyramid Principle}~\cite{minto2009pyramid}~:
          la décomposition d'un problème en éléments exclusifs et
          exhaustifs.
    \item A0, \textit{Vers l'auditologie}~\cite{minka2026auditologie}~: le
          gabarit canonique dans le programme de recherche.
\end{itemize}
\end{boxinfo}

\begin{boxinfo}
\textbf{Transition.} La décomposition dit quoi vérifier et comment. Il
reste à aller chercher les réponses et les preuves sur le terrain. Le
Chapitre~\ref{chap:preuve} présente l'instrument \OUI{}, \NON{} ou \NEK{},
les entretiens, les tests, l'échantillonnage et les papiers de travail.
\end{boxinfo}
